质量 \(m\) 的小珠在水平无摩擦平面上沿由中心控制的细直槽滑动，珠到中心的距离由理想执行器强制为 \(r(t)=r_0(1+\eta t)>0\)，其中 \(r_0>0,\eta>0\)。槽能自由绕固定中心旋转，且槽自身的转动惯量可忽略。初始珠角速度为 \(\omega_0>0\)，初态径向速度为 \(r_0\eta\)。忽略空气阻力和轴摩擦。
(1) 证明珠关于中心的角动量守恒，求 \(\dot\phi(t)\) 与转角函数。
(2) 求切向约束力和执行器必须施加的径向力，分别指出它们的作功率。
(3) 求从时刻零到时刻 \(t\) 珠的动能变化和执行器总作功，检查机械功收支。
(4) 计算角速度趋零时的渐近规律，并解释执行器反向拉回时能量回收条件。
